\newcommand{\DoNotLoadEpstopdf}{}
\pdfinfoomitdate=1
\pdftrailerid{}
\pdfsuppressptexinfo=-1
\documentclass[11pt,letterpaper]{article}
\usepackage[T1]{fontenc}
\usepackage{lmodern}
\usepackage{amsmath,amssymb,amsthm,mathtools}
\usepackage[margin=1in]{geometry}
\usepackage{booktabs,array,microtype}
\usepackage[hidelinks]{hyperref}
\hypersetup{pdftitle={Report 167 Asymptotics of assemblies of nested cycles},pdfauthor={},pdfsubject={Fixed order marked saddle expansion for A308338 and A392471}}
\setlength{\parskip}{4pt}
\setlength{\parindent}{1.2em}
\setlength{\emergencystretch}{3em}
\setcounter{tocdepth}{1}
\numberwithin{equation}{section}
\newtheorem{theorem}{Theorem}[section]
\newtheorem{lemma}[theorem]{Lemma}
\newtheorem{proposition}[theorem]{Proposition}
\newtheorem{corollary}[theorem]{Corollary}
\theoremstyle{remark}\newtheorem{remark}[theorem]{Remark}
\newcommand{\dd}{\,\mathrm d}
\newcommand{\CC}{\mathbb C}
\newcommand{\NN}{\mathbb N}
\newcommand{\EE}{\mathbb E}
\newcommand{\Var}{\operatorname{Var}}
\newcommand{\Li}{\operatorname{Li}}
\newcommand{\dist}{\operatorname{dist}}
% Added by the ProveIt write phase (batch 108, 6 October 2026): dated notes
% marked [write] and a macro for file, label and declaration names.
\newcommand{\file}[1]{\nolinkurl{#1}}
\expandafter\def\expandafter\UrlBreaks\expandafter{\UrlBreaks\do\-\do\_\do\.}
\newenvironment{writenote}[1][]{\par\smallskip\begingroup\small
 \noindent\textbf{[write]}\if\relax\detokenize{#1}\relax\else~\textbf{#1.}\fi\ }%
 {\par\endgroup\smallskip}
\title{Asymptotics of assemblies of nested cycles\\\large Exact counts and fixed order marked saddle expansions}
\author{Report 167}
\date{3 October 2026}
\begin{document}
\maketitle
\begin{abstract}
Let $P(z)=\prod_{m\geq1}(1+z^m/m)$ and
$F(z,u)=\exp\{u(P(z)-1)\}$. The labeled count
$a_n=n![z^n]F(z,1)$ is OEIS A308338, and the coefficient of $u^k$ counts
objects with $k$ nonempty nested-cycle components, giving A392471.
With $c=e^{-\gamma}$, we prove
\[
 a_n\sim \frac{e^{c(1/2-\log2)-1}c^{1/4-c/2}}{2\sqrt\pi}
             n!n^{c/2-3/4}e^{2\sqrt{cn}}.
\]
More generally, every fixed order has a complex-uniform expansion in
$n^{-1/2}$ with polynomial coefficients in $\log n$; we give the first
two corrections explicitly. A whole-circle estimate controls all
secondary arcs, and a Gaussian Taylor lemma gives the stated sharp
logarithmic remainder without crossing a natural boundary.

For the component count, holomorphic normalization justifies
differentiation and gives explicit logarithmic and constant corrections
to the mean and variance. In particular the mean divided by $\sqrt n$
tends to $e^{-\gamma/2}$. We recover the established Gaussian regime
with variance equivalent $\sqrt{cn}/2$ and prove a rounding-safe
asymptotic inverse for the integer count sequence.

The component asymptotics and generic expansive-assembly limit laws
are prior results. The contribution developed here is the explicit
sequence-specific marked expansion and its refinements. The companion
contains bounded exact integer and symbolic checks. No worldwide
priority claim, complete beyond-all-orders transseries, or certified
numerical asymptotic error bound is asserted.
\end{abstract}
\begin{writenote}[Status and trust boundary, added 6 October 2026, batch~108]
This is bundle Report~167 of an external AI-assisted research session, filed
in ProveIt as the single-source research report
\file{a308338-nested-cycle-assemblies} (provenance, the sources as the write
read them, relation to the repository, the collected non-claims and the
reading conventions in Section~\ref{ncy:sec:provenance}). It is unrefereed
and not formalized; its author line reads ``Report 167'' and its PDF author
field is empty. \emph{What rests on what.} Everything analytic rests on the
component estimate~\eqref{ncy:eq:prior}, credited to Greene and Knuth through
Flajolet, Fusy, Gourdon, Panario and Pouyanne and cited, not reproved. Given
it, the text proves the fixed-order expansion
(Theorem~\ref{ncy:thm:main}), the moment expansions
(Theorem~\ref{ncy:thm:moments}), the Gaussian limit
(Corollary~\ref{ncy:cor:clt}, a known regime re-proved with this
normalization) and the two-ceiling inverse (Theorem~\ref{ncy:thm:inverse});
no constant or onset is made effective and every decimal is uncertified.
\emph{Two statements of the literature are put on record, with proofs.}
(i)~Riedel's January~2026 conjecture $\mathrm E[X]\sim N\sqrt n$ for the
number of components is \emph{proved} by~\eqref{ncy:eq:meanlimit}, with
$N=e^{-\gamma/2}=0.7493\ldots$; this excludes the alternative he raised, a
power $n^M$ with $M$ ``close to but not equal to $1/2$'', and his numerical
$N\approx0.7719351399$ is not the constant but agrees to nine decimals with
the finite ratio $\EE K_{4500}/\sqrt{4500}$ (Remark~\ref{ncy:rem:riedel}). The source quotes the
estimate without saying that it settles the conjecture. (ii)~The general-$C$
covariance display of Erlihson and Granovsky (preprint~(4.54),
published~(4.18), both read) is negative at $p=1$, $C=3$, $u=1$ and
contradicts Corollary~\ref{ncy:cor:clt} in this report's own model; the write
gives the corrected display and locates the dropped factor~$C$
(Remark~\ref{ncy:rem:eg}). The source had already found the negative value.
The write also adds Remark~\ref{ncy:rem:transseries} (the inversion against
the transseries volume), notes on the OEIS entries as of 6~October~2026, and one question
(Section~\ref{ncy:sec:questions}). Nothing was submitted to the OEIS or sent to
any author.
\end{writenote}
\newpage
\tableofcontents
\medskip
\noindent\textbf{Scope and reading guide.}
Sections~\ref{ncy:sec:exact}--\ref{ncy:sec:main} specify the object and results.
Sections~\ref{ncy:sec:sector}--\ref{ncy:sec:gaussian} prove the fixed-order
coefficient theorem. Sections~\ref{ncy:sec:corrections}--\ref{ncy:sec:inverse}
give corrections, probability laws, and inversion.
Section~\ref{ncy:sec:computation} separates exact verification from numerical
diagnostics. Section~\ref{ncy:sec:sources} records source coverage and
limits. The report is self-contained after the explicitly credited
component coefficient estimate~\eqref{ncy:eq:prior}.
\newpage

\section{Exact objects and established input}\label{ncy:sec:exact}
A labeled cycle on an $m$-element set has $(m-1)!$ possible cyclic
orders and exponential generating function $z^m/m$. Choosing at most
one cycle of each length therefore gives
\begin{equation}\label{ncy:eq:P}
 P(z)=\prod_{m\geq1}\left(1+\frac{z^m}{m}\right)
     =\sum_{m\geq0}B_m\frac{z^m}{m!},\qquad p_m=\frac{B_m}{m!}.
\end{equation}
The integers $B_m$ form A007838. A nonempty collection of cycles of
distinct lengths has a unique decreasing-length nesting order. It is
one nested component in Riedel's interpretation; there is no extra
ordering or geometric-embedding choice. A set of such nonempty labeled
components has marked exponential generating function
\begin{equation}\label{ncy:eq:F}
 F(z,u)=\exp\{u(P(z)-1)\},\qquad f_n(u)=[z^n]F(z,u).
\end{equation}
Thus
\begin{equation}\label{ncy:eq:counts}
 T(n,k)=\frac{n!}{k!}[z^n](P(z)-1)^k,
 \qquad a_n=n!f_n(1)=\sum_{k=0}^nT(n,k).
\end{equation}
The $a_n$ are A308338, with $a_0=1$, and the published triangle A392471
starts at $n=1$, $1\leq k\leq n$. We use the natural empty extension
$T(0,0)=1$ and $T(n,0)=0$ for $n>0$ only for recurrences. In particular
$T(n,1)=B_n$ for $n\geq1$.

Ilya Gutkovskiy entered A308338 and its defining EGF on 20 May 2019
\cite{oeis308338}. Riedel's January 2026 note supplies the nested-cycle
interpretation, the marked triangle, and recurrences
\cite{riedel,oeis392471}. The component EGF is much older. The prior
estimate we use is
\begin{equation}\label{ncy:eq:prior}
 p_m=c+\frac{c}{m}+O\!\left(\frac{\log m}{m^2}\right),
 \qquad c=e^{-\gamma},\qquad m\longrightarrow\infty,
\end{equation}
where $\gamma$ is Euler's constant. Flajolet, Fusy, Gourdon, Panario,
and Pouyanne \cite[Section 3, equation (22)]{hybrid} attribute this
estimate to Greene and Knuth. Their Proposition 1 gives the full
component coefficient expansion, including periodic terms from roots
of unity. We use~\eqref{ncy:eq:prior} as a credited theorem; their theorem
for $P$ is not itself a coefficient theorem for the outer exponential
$F$.

For a component spectrum $(j_1,\ldots,j_n)$ with $\sum m j_m=n$, the
number of objects is
\[
 n!\prod_{m=1}^n\frac{p_m^{j_m}}{j_m!}.
\]
After division by $a_n$, this is the Gibbs assembly law with weights
$p_m\sim c$, so its expansive exponent is $1$. The weights are $p_m$,
not $B_m$ or $p_m/m$. The square-root component law and normal
fluctuations lie within established expansive-assembly theory.
Erlihson and Granovsky \cite[Corollary 4.5]{eg2008} explicitly attribute
the total-component CLT to their 2004 work~\cite{eg2004}. We derive the
precise normalization and lower-order terms directly below.

\subsection{Provenance, sources and reading conventions}\label{ncy:sec:provenance}
\begin{writenote}[Added 6 October 2026, batch~108]
\emph{Provenance.} This report is Report~167 of the session bundle that
ProveIt's \file{docs/incoming} intake processed as batch~108: the archive
\file{Nested_Cycle_Assemblies_Asymptotics_and_Inverses_Source.zip}
(938{,}156 bytes, SHA-256 \file{ef59faff7aba...eaebc0b63}, 21 files at the
archive root and in its \file{companion/} and \file{data/} folders), arrival
commit \file{60f54ea06}, placed in \file{602e5bd0f}. The text is the
delivered \file{Report167.tex} (908 lines, dated 3~October 2026), shipped as
\file{article.tex}. The delivered 15-page PDF and the checksum manifest
\file{SHA256SUMS} (20 entries, verified at placement and again at the write)
are not shipped, and the delivered README is replaced by the report README;
all three are retrievable from the arrival commit. The package records no
ProveIt commit (it cites nothing in the repository), no ``prepared for
private review'' line and no personal data; the word ``private'' occurs only
in the delivered README's ``private fresh format'' (a \TeX\ format file). It
is a single-source report, so no merge choice was made. The write prefixed
the 62 delivered labels with \file{ncy:} and updated the 60 references to
them; it added this subsection, the notes marked [write],
Remarks~\ref{ncy:rem:riedel}, \ref{ncy:rem:transseries} and~\ref{ncy:rem:eg},
labels on the items of Section~\ref{ncy:sec:questions} and one further item
there. No statement, proof or number of the source was changed: the added
remarks are the last statements of their sections, the added displays are
unnumbered, and the added item follows the last delivered one, so every
theorem, section and equation of the source keeps its delivered number.

\emph{The sources, as the write read them} (6~October 2026).
\begin{itemize}\raggedright
\item OEIS A308338~\cite{oeis308338}, live revision~\#21 (5~October 2026);
 the source froze revision~\#17 (15~January 2026). Revisions \#18--\#21, by
 Alois P. Heinz on 5~October 2026, add a Maple program and a b-file for
 $0\le n\le445$; the 22 data terms are unchanged. Riedel's comment
 identifying the terms as nested cycle partitions and his formula
 $a(n)=\sum_kT(n,k)$ date from 14~January 2026. The entry states no
 asymptotic and no conjecture.
\item OEIS A392471~\cite{oeis392471}, live revision~\#42 (5~October 2026);
 the source froze revision~\#37 (25~January 2026). On 5~October 2026 Heinz
 inserted the column $k=0$ (the entry's \texttt{\%E} line reads ``Column k=0
 inserted by Alois P. Heinz, Oct 05 2026''), so the entry now has offset~$0$
 and rows from $n=0$, and he added a Maple program and a b-file of rows
 $0\le n\le150$. The delivered sentence of Section~\ref{ncy:sec:exact} that
 the published triangle ``starts at $n=1$, $1\leq k\leq n$'' describes
 revision~\#37; the entry now uses the same extension $T(0,0)=1$, $T(n,0)=0$
 ($n>0$) as this report. Its link to Riedel's original note is titled
 ``Nested Cycle Partitions: definitions and basic recurrences''; the PDF
 behind it is titled ``Nested Cycle Partitions: a conjecture''. The entry
 itself contains no conjecture.
\item OEIS A007838~\cite{oeis7838}, revision~\#74 (13~March 2026), unchanged;
 its formula section states~\eqref{ncy:eq:prior} in the form
 $a(n)\sim n!(e^{-g}+e^{-g}/n+O((\log n)/n^2))$ and its reference section
 cites Greene and Knuth.
\item The three b-files agree with \file{data/exact_data_100.json} for
 $0\le n\le100$ (A308338, A007838) and for all rows $0\le n\le100$ of the
 triangle (the write's comparison).
\item Riedel's note~\cite{riedel} in both versions: the revised one
 (author-hosted, page footers dated 24.01.26) and the original linked from
 A392471 (footers dated 13.01.26), six pages each; the paragraph
 ``The conjecture'' on p.~3 is the same in both
 (Remark~\ref{ncy:rem:riedel}).
\item Erlihson and Granovsky~\cite{eg2008}: the preprint
 arXiv:math/0507343v3 (19~June 2007), Section~4 from~(4.37) to~(4.55) and
 Section~5 around~(5.121), and the published text (numdam), pp.~925--927
 and~936--938 (Remark~\ref{ncy:rem:eg}).
\item The transseries volume
 \path{Analysis/Transseries/docs/series-and-transseries/Transseries_And_Inversion/transseries_and_inversion.tex}:
 \file{p0:def:core}, \file{p0:prop:factorial-core},
 \file{p0:def:three-inverses} and \file{p0:thm:staircase}, read for
 Remark~\ref{ncy:rem:transseries}.
\item Not read by the write: the hybrid paper~\cite{hybrid} (the source read
 its Section~3 and Proposition~1), the Greene--Knuth book and the 2004 paper
 of Erlihson and Granovsky~\cite{eg2004} (the source read neither).
 \file{data/SOURCE_PROVENANCE.json} records what the source inspected.
\end{itemize}

\emph{What was checked.} The batch-108 dossier read the manuscript in full
and found no false claim. It checked~\eqref{ncy:eq:prior} at $m=100$
($p_{100}=0.56678641$ against $c+c/100=0.56707408$, a difference within
$\log m/m^2$), the counts, the exact mean and the exact variance at $n=100$
against~\eqref{ncy:eq:equivalent} and Theorem~\ref{ncy:thm:moments}, read
Riedel's conjecture paragraph, verified the literal substitution in the
Erlihson--Granovsky preprint, and replayed the companion on copies (all four
frozen outputs reproduced). At the write the companion was rerun from a fresh
extraction of the archive and from the shipped files (note at the end of
Section~\ref{ncy:sec:computation}), and $f_n(1)$, $\EE K_n$ and $\Var K_n$
were evaluated in double precision for $n\le40000$ by recurrences with
positive terms only ($p_m$ from the cycle product
of~\eqref{ncy:eq:integercycle} divided by $m!$, $f_n(1)$ from the second
recurrence of~\eqref{ncy:eq:rational}, and the moments from
$[z^n]F(z,1)(P(z)-1)$ and $[z^n]F(z,1)(P(z)-1)^2$), agreeing with the exact
rational values at
$n=100$ ($\EE K_{100}=8.219543$, $\Var K_{100}=4.016577$). The residuals are
consistent with Theorem~\ref{ncy:thm:main} at $J=1$ and with
Theorem~\ref{ncy:thm:moments} including $M_1$ and $V_1$: at $n=40000$,
$\sqrt n\,(a_n/(Cn!n^{c/2-3/4}e^{2\sqrt{cn}})-1)=9.682$ against
$R_1(\log n;1)=9.344$, $\sqrt n$ times the residual of~\eqref{ncy:eq:mean} is
$12.92$ against $M_1(\log n)=12.69$, and $\sqrt n$ times the residual
of~\eqref{ncy:eq:variance} is $16.97$ against $V_1(\log n)=16.53$; at
$n=1000$ the same pairs are $5.43/4.71$, $6.74/6.13$ and $8.38/7.27$, so the
gaps shrink as the next order predicts. These are observations, not
certificates.

\emph{Relation to the repository.} Before batch~108 no file of the
repository named A308338, A392471 or A007838, nested cycle partitions or
Riedel's note, and none does outside this report and the collection
catalogue. No statement of this report is formalized in Lean or Rocq, and its
place in the collection confers no formal status. The neighbouring reports
\file{a064856-stirling-catalan-transforms} (Part~II: cycle-weighted
permutation mixtures) and \file{a301746-divisor-weighted-asymptotics},
\file{a022629-distinct-partition-norms}, \file{a271619-strict-twice-partitions}
(Gibbs partitions after Granovsky, Stark and Erlihson) treat other models;
none uses the Erlihson--Granovsky display corrected in
Remark~\ref{ncy:rem:eg}, and none needs a reciprocal note. The inversion of
Section~\ref{ncy:sec:inverse} is compared with the transseries volume in
Remark~\ref{ncy:rem:transseries}.
\end{writenote}
\begin{writenote}[What is not claimed, collected from the source]
The expansion of Theorem~\ref{ncy:thm:main} holds at every fixed order only:
no convergence, optimal truncation, Stokes phenomenon or complete
beyond-all-orders description is claimed, and constants need not be uniform
when the order grows. The decimal value of $C$ is an uncertified floating
evaluation; no interval enclosure of $C$ is supplied, and no implied constant
or onset (in particular no certified $M_J$ or threshold in
Theorem~\ref{ncy:thm:inverse}) is made effective. The calculator
\eqref{ncy:eq:calculator} is formal: its contour is not identified with a
Bessel contour and no convergence is claimed. The square-root law, the exact
limit~\eqref{ncy:eq:meanlimit} and the normal fluctuations belong to the
established expansive-assembly regime; Corollary~\ref{ncy:cor:clt} is ``not a
new universality assertion'', and no local limit theorem or distribution
distance bound is inferred. The finite and symbolic checks prove no
asymptotic statement; the $B$, $B_3$ checks verify the arithmetic pieces, not
the infinite-tail summations. Source coverage is bounded: no worldwide
priority is claimed, ``No claim is made that no other work contains
overlapping results'', the Greene--Knuth book and the 2004 paper were not
inspected, and no OEIS entry was amended. Byte identity of rebuilds is
toolchain-specific. The write adds: the table and the identification of
$n=4500$ in Remark~\ref{ncy:rem:riedel} are floating evaluations (checked
against exact and 256-bit fixed-point values) and an inference about how
Riedel's number arose, not a statement of his; the corrected display of
Remark~\ref{ncy:rem:eg} is derived from the conditioning computation, not
from a re-audit of the whole proof of Erlihson and Granovsky.
\end{writenote}
\medskip\noindent\begin{minipage}{\textwidth}
\emph{Reading conventions} (letters the source reuses, and letters of the
two outside sources discussed by the write, with the tempting false
readings; no symbol of the source was renamed; $\varrho$ and $M_\varrho(n)$
are the write's).\par
\smallskip
\begingroup\small\renewcommand{\arraystretch}{1.1}
\noindent\begin{tabular}{@{}>{\raggedright\arraybackslash}p{1.0in}>{\raggedright\arraybackslash}p{2.9in}>{\raggedright\arraybackslash}p{2.2in}@{}}
\toprule
Symbol & Meaning here (where) & Not to be read as\\
\midrule
$B_m$, $p_m$ & component counts (A007838) and $B_m/m!$ & the constants $B$, $B_3$ of~\eqref{ncy:eq:logP}; the calculator $\mathcal B_\beta$; Erlihson--Granovsky's $b_r$ and exponent $p$\\
$c$, $A$, $D$, $E$ & $e^{-\gamma}$; constants of~\eqref{ncy:eq:C}, \eqref{ncy:eq:QZ} & OEIS A-numbers; $\EE$; the residual $E_{n,J}$\\
$C(u)$, $C$ & amplitude~\eqref{ncy:eq:C}, \eqref{ncy:eq:equivalent} & $C_0$, $C_{j,k,m}$ (proof constants); Erlihson--Granovsky's $C$ in $a_k\sim Ck^{p-1}$\\
$u$, $v=cu$ & marking variable; $cu$ & Erlihson--Granovsky's scaled size $u$; $v$ of \file{p0:prop:factorial-core}\\
$L$ & $\log n$ in $R_j(L;u)$, $M_1(L)$, $V_1(L)$; $\log y$ in Section~\ref{ncy:sec:inverse} & $L_n(u)$ (leading term); the target $L$ of \file{p0:prop:factorial-core}\\
$Q(X)$, $Q_j$, $X_n$ & polynomials of~\eqref{ncy:eq:QZ}, \eqref{ncy:eq:Qj}; contour bound & Riedel's $Q(z,u)=F(z,u)$ and $Q_k(z)$; Riedel's $X=K_n$\\
$K_n$, $K$, $K_0$ & component count; constant in the proof of Theorem~\ref{ncy:thm:inverse}; $\log(C\sqrt{2\pi})$ & \\
$H_{j,k}$, $H_j$ & Taylor factor~\eqref{ncy:eq:H}; harmonic numbers & \\
$\ell$, $\ell_c$, $\ell_r$ & $\log s$; $(\log c-L)/2$; coefficients of $\log P$ (Section~\ref{ncy:sec:computation}) & \\
$\beta$ & $c/2-3/4$ (Section~\ref{ncy:sec:inverse}) & the free exponent in $\mathcal B_\beta$\\
$\kappa$, $\kappa'$, $h$, $d$, $\delta$ & loss constants, window $a^{3/2}\log n$ (Section~\ref{ncy:sec:circle}); bound and $n^{-1/4}$ in Lemma~\ref{ncy:lem:gaussian} & the volume's $\kappa_{\mathrm{vol}}$, $d_{\mathrm{vol}}$, $h$; Riedel's $\kappa_n$; Erlihson--Granovsky's $\delta=1/r_N$\\
$p$, $q$ & $2J+2$, $(J+1)/2$ (Section~\ref{ncy:sec:inverse}); $q=Q(\ell_c)$ in~\eqref{ncy:eq:explicitM1} & $p_m$; Erlihson--Granovsky's exponent $p$ and number $q$ of strata\\
$N(y)$, $M_J$, $M_1$ & first passage; constant of Theorem~\ref{ncy:thm:inverse}; mean correction & Riedel's constant $N$ and exponent $M$; Erlihson--Granovsky's size $N$; the write's $M_\varrho(n)$\\
$T(n,k)$, $t$ & A392471 (with $T(n,0)$); $z=e^{-t}$ & \\
\bottomrule
\end{tabular}
\endgroup
\end{minipage}

\section{The marked coefficient theorem}\label{ncy:sec:main}
Fix a sufficiently small open complex disc about $u=1$. Choose
$\log(cu)$ and $\sqrt{cu}$ to be real on positive real $u$, and set
\begin{equation}\label{ncy:eq:C}
 A=\frac12-\log2,\qquad
 C(u)=\frac{e^{u(cA-1)}(cu)^{1/4-cu/2}}{2\sqrt\pi}.
\end{equation}
All statements about fixed order mean that the order is fixed as
$n\to\infty$; constants need not be uniform when that order grows.

\begin{theorem}[Uniform fixed-order expansion]\label{ncy:thm:main}
For every fixed integer $J\geq0$ there are polynomials $R_j(L;u)$ of
degree at most $2j$ in $L$, with coefficients holomorphic near $1$, such
that
\begin{align}\label{ncy:eq:main}
 f_n(u)&=C(u)n^{cu/2-3/4}e^{2\sqrt{cun}}
 \left\{1+\sum_{j=1}^J n^{-j/2}R_j(\log n;u)
       +E_{n,J}(u)\right\},\\
 E_{n,J}(u)&=O\!\left(n^{-(J+1)/2}(1+\log n)^{2J+2}\right).
 \notag
\end{align}
The estimate is uniform on every sufficiently small fixed closed disc
about $1$. The exact normalized residual $E_{n,J}$ is holomorphic on
an open disc containing any smaller disc on which derivatives are
taken, and each fixed derivative obeys the same order bound there.
\end{theorem}

Taking $u=1$, $J=0$ gives the equivalent
\begin{equation}\label{ncy:eq:equivalent}
 a_n\sim C n!n^{c/2-3/4}e^{2\sqrt{cn}},\qquad
 C=\frac{e^{c(1/2-\log2)-1}c^{1/4-c/2}}{2\sqrt\pi}.
\end{equation}
For orientation only,
$C\approx0.0947779253874791584250013230745944160$; this is an
uncertified floating evaluation of the exact expression. The logarithmic
term of $P$ supplies the power $n^{c/2}$ in~\eqref{ncy:eq:equivalent}.
Discarding that term would give an incorrect equivalent.

The theorem is an all-fixed-orders algebraic expansion with logarithmic
polynomials. It makes no claim about convergence, optimal truncation,
Stokes phenomena, or a complete description of terms beyond all orders.

\section{Local component expansion in a right half plane}\label{ncy:sec:sector}
We first prove the needed local expansion without continuing $P$
through its natural boundary. For $\Re t>0$, isolating the factor $m=1$
gives the absolutely convergent identity
\begin{align}\label{ncy:eq:polylog}
 \log P(e^{-t})={}&\log(1+e^{-t})+\Li_1(e^{-t})-e^{-t}\\
 &+\sum_{k\geq2}\frac{(-1)^{k+1}}{k}
       \{\Li_k(e^{-kt})-e^{-kt}\}.\notag
\end{align}
Fix a desired order $H$ and choose an integer $K>H+2$. The tail $k>K$
can be written as a double sum over $k>K$, $m\geq2$. For
$0\leq h\leq H+1$, its $h$th derivative in $t$ is bounded absolutely
and uniformly on the closed right half plane by
\begin{equation}\label{ncy:eq:tail}
 \sum_{k>K}k^{h-1}\sum_{m\geq2}m^{h-k}<\infty.
\end{equation}
Indeed $k-h>1$ and the inner sum decreases exponentially as $k$ grows.
Termwise differentiation up to that order is consequently valid. The
integral remainder along the segment from $0$ to $t$ gives a Taylor
polynomial for the tail with error $O(|t|^{H+1})$, uniformly in the
closed right half plane near $0$.

Each retained polylogarithm has the classical finite-order expansion
\begin{equation}\label{ncy:eq:li}
 \Li_k(e^{-x})=
 \frac{(-x)^{k-1}}{(k-1)!}(H_{k-1}-\log x)
 +\sum_{\substack{r\geq0\\r\ne k-1}}
       \zeta(k-r)\frac{(-x)^r}{r!},
\end{equation}
where $H_j=\sum_{m=1}^j1/m$ and $H_0=0$. For a fixed finite set of
$k$, truncations and their remainders are uniform in every closed
subsector of $\Re t>0$. Combining them with~\eqref{ncy:eq:tail} gives each
fixed order of $\log P(e^{-t})$. At every positive power $t^r$, the
coefficient is affine in $\log t$. Its logarithmic coefficient comes
only from $k=r+1$ in~\eqref{ncy:eq:polylog} and is
\[
 -\frac{(r+1)^{r-1}}{r!}\qquad(r\geq1).
\]
The constant after removing $-\log t$ is $-\gamma$, by the convergent
Euler product with the harmonic term subtracted. Explicitly,
\begin{align}\label{ncy:eq:logP}
 \log P(e^{-t})={}&-\log t-\gamma+t(-\log t+A)
       +t^2\left(-\frac32\log t+B\right)\notag\\
 &+t^3\left(-\frac83\log t+B_3\right)
       +O\bigl(|t|^4(1+|\log t|)\bigr),\\
 B={}&\frac{41}{24}-\frac{\zeta(2)}2-\frac32\log3,\qquad
 B_3=\frac{49}{12}-\frac{\pi^2}{4}
             +\frac{\zeta(3)}3-\frac83\log4.\notag
\end{align}

For completeness, the $t^2$ constant contributions from the isolated
prefactor, $k=2$, $k=3$, and $k\geq4$ are, respectively,
\[
 -\frac5{12},\qquad \frac32,\qquad
 \frac34-\frac32\log3,\qquad -\frac{\zeta(2)}2-\frac18.
\]
The last follows by summing a rational series:
\begin{align*}
 \frac12\sum_{k\geq4}(-1)^{k+1}k(\zeta(k-2)-1)
 &=-\frac12\sum_{m\geq2}\frac{4m+3}{m(m+1)^2}\\
 &=-\frac{\zeta(2)}2-\frac18.
\end{align*}
For $B_3$, the corresponding contributions (now separating $k=4$)
are $1/6$, $-13/18$, $9/4$, $20/9-(8/3)\log4$, and
\[
 -\frac16\sum_{m\geq2}\frac{25m^2+39m+16}{m(m+1)^3}
   =\frac16-\frac32\zeta(2)+\frac13\zeta(3).
\]
These calculations supply the constants rather than attributing them
to a short displayed formula in the component source.

Exponentiation to any prescribed fixed order yields
\begin{equation}\label{ncy:eq:localP}
 P(e^{-t})=\frac ct+c(-\log t+A)
     +\sum_{j=1}^{M}t^jU_j(\log t)
     +O\bigl(|t|^{M+1}(1+|\log t|)^{M+2}\bigr),
\end{equation}
where $\deg U_j\leq j+1$. To obtain the stated error, first expand
$\log P$ through degree $M+1$, with error
$O(|t|^{M+2}(1+|\log t|))$, and then exponentiate. The largest omitted
logarithmic degree comes from $M+2$ copies of $t\log t$.
Define
\begin{align}\label{ncy:eq:QZ}
 D&=\log2-2,\qquad E=B+A^2/2,\qquad Q(X)=X^2/2+DX+E,\notag\\
 Z(X)&=\frac{(-X+A)^3}{6}+(-X+A)\left(-\frac32X+B\right)
                  -\frac83X+B_3.
\end{align}
Then $U_1=cQ$ and $U_2=cZ$. All these are sector expansions inside the
right half plane. No analytic continuation of $P$ across $|z|=1$ has
been used.

\section{The complete circle and complex saddle}\label{ncy:sec:circle}
The global step comes from the prior coefficient estimate rather than
from a claim that the positive real singularity is the only singularity.
From~\eqref{ncy:eq:prior}, absolute summation of
$|p_m-c|e^{-am}$, including the bounded $m=0$ adjustment, gives uniformly
for $0<a<1$, $|y|\leq\pi$,
\begin{equation}\label{ncy:eq:globalP}
 P(e^{-a-iy})=\frac{c}{1-e^{-a-iy}}+O(1+\log(1/a))
             =\frac{c}{a+iy}+O(1+\log(1/a)).
\end{equation}
For the second equality, $(1-e^{-t})^{-1}-t^{-1}$ is bounded on
$0\leq\Re t\leq1$, $|\Im t|\leq\pi$ after filling its removable
value at $0$. Thus~\eqref{ncy:eq:globalP} includes every secondary
root-of-unity arc.

Put $v=cu$ and $s=\sqrt{v/n}=a+ib$, where $a>0$. Coefficient extraction
on $|z|=e^{-a}$, with $z=e^{-t}$, gives
\begin{equation}\label{ncy:eq:cauchy}
 f_n(u)=\frac1{2\pi i}\int_{a-i\pi}^{a+i\pi}
        \exp\{u(P(e^{-t})-1)+nt\}\dd t.
\end{equation}
The sign from $\dd z/z=-\dd t$ is canceled by reversal of the circle
parameter. The principal phase is $\Phi(t)=v/t+nt$, with saddle $s$.
Since $v=ns^2$, direct algebra gives the exact loss identity
\begin{equation}\label{ncy:eq:loss}
 \Re\Phi(a+iy)=2na-na\frac{(y-b)^2}{a^2+y^2}.
\end{equation}
Uniformly on the small marking disc, $a\asymp n^{-1/2}$ and $|b|/a$
is bounded. Take $h=a^{3/2}\log n$. Outside $|y-b|<h$, the loss in
\eqref{ncy:eq:loss} is at least $\kappa(\log n)^2$ for a positive
constant $\kappa$. To see this, when $|y|\leq C_0a$ the denominator
is $O(a^2)$ and $na h^2/a^2\asymp(\log n)^2$; when
$|y|>C_0a$ with $C_0$ sufficiently large, the fraction is bounded
below, and the loss is of order $\sqrt n$.

The error from~\eqref{ncy:eq:globalP} costs only $O(\log n)$ in the real
exponent, with $u$ bounded. The entire off-saddle contribution is
therefore
\begin{equation}\label{ncy:eq:off}
 O\!\left(e^{2na-\kappa'(\log n)^2}\right)
\end{equation}
for a uniform $\kappa'>0$. The modulus of the leading term in
\eqref{ncy:eq:main} is $e^{2na}$ times a quantity bounded above and below
by fixed powers of $n$. Hence~\eqref{ncy:eq:off}, relatively to that term,
is smaller than every fixed inverse power of $n$. This proves the
needed suppression for the whole circle, including the natural-boundary
points, rather than relying only on positivity or aperiodicity.

\section{The local Gaussian expansion and its remainder}\label{ncy:sec:gaussian}
On $|t-s|\leq h$, $t$ remains in a fixed right-half-plane subsector
and $|t|\asymp n^{-1/2}$. Equation~\eqref{ncy:eq:localP} gives, for any
fixed $J$,
\begin{align}\label{ncy:eq:amplitude}
 e^{u(P(e^{-t})-1)}={}&e^{u(cA-1)}t^{-v}e^{v/t}\\
 &\quad\times\left\{\sum_{j=0}^Jt^jQ_j(\log t;u)
   +O\bigl(|t|^{J+1}(1+|\log t|)^{2J+2}\bigr)\right\}.\notag
\end{align}
Here $Q_0=1$, and all $Q_j$ are defined by the formal identity
\begin{equation}\label{ncy:eq:Qj}
 \exp\left(u\sum_{j\geq1}t^jU_j(X)\right)
             =\sum_{j\geq0}t^jQ_j(X;u).
\end{equation}
A product of $r$ positive-degree factors with total $t$-degree $j$
has logarithmic degree at most $j+r\leq2j$; therefore
$\deg_XQ_j\leq2j$. Truncating the ordinary exponential gives the error
in~\eqref{ncy:eq:amplitude}, since $t(1+|\log t|)^2\to0$ on the segment.
The component error in~\eqref{ncy:eq:localP} is no larger. In particular,
\begin{equation}\label{ncy:eq:Q12}
 Q_1=vQ,\qquad Q_2=vZ+\frac{v^2}{2}Q^2.
\end{equation}

The next lemma records the integrated Taylor remainder carefully;
replacing powers of the integration variable by endpoint bounds would
lose unnecessary powers of $\log n$.

\begin{lemma}[Uniform Gaussian Taylor remainder]\label{ncy:lem:gaussian}
Set $\delta=n^{-1/4}$, $\lambda=\sqrt{vn}$, and use the actual
vertical contour parametrization
\[
 t=s+in^{-3/4}x=s(1+i\delta x/\sqrt v),\qquad
 |x|\leq X_n=(\Re\sqrt v)^{3/2}\log n.
\]
For fixed nonnegative integers $j,k,S$, the integral of the monomial
$t^{j-v}(\log t)^k$ times $e^{nt+v/t}$ has an expansion, after removing
its common leading scale, in powers $n^{-r/2}$, $0\leq r\leq S$.
Its relative-to-that-scale remainder is
$O(n^{-(S+1)/2}(1+|\log s|)^k)$, uniformly in $u$ on a sufficiently
small fixed closed disc. Each coefficient is polynomial in $\log s$
of degree at most $k$ and holomorphic in $u$.
\end{lemma}
\begin{proof}
The exact phase, with no truncated exponential yet, is
\begin{equation}\label{ncy:eq:exactphase}
 nt+\frac vt=2\lambda-\frac{x^2}{\sqrt v}
             +\frac{i\delta x^3}{v(1+i\delta x/\sqrt v)}.
\end{equation}
Write $\ell=\log s=(\log v-\log n)/2$ and remove $s^{j-v}$ from the
monomial. The residual analytic factor multiplying the Gaussian is
\begin{align}\label{ncy:eq:H}
 H_{j,k}(\delta,x;v,\ell)={}&(1+i\delta x/\sqrt v)^{j-v}
       [\ell+\log(1+i\delta x/\sqrt v)]^k\notag\\
 &\times\exp\left(\frac{i\delta x^3}{v(1+i\delta x/\sqrt v)}\right).
\end{align}
In a Taylor expansion in $\delta$, hold $\ell$ fixed; substitute its
stated value afterwards. Let $d>0$ be a uniform lower bound for
$\Re(v^{-1/2})$ on the marking disc. For every fixed derivative order
$m$ and all sufficiently large $n$, one has, for $0\leq\xi\leq\delta$
and $|x|\leq X_n$,
\begin{equation}\label{ncy:eq:derivbound}
 |\partial_\delta^mH_{j,k}(\xi,x;v,\ell)|
 \leq C_{j,k,m}(1+|\ell|)^k
            P_{j,k,m}(|x|)e^{dx^2/2},
\end{equation}
with a fixed polynomial $P_{j,k,m}$ having nonnegative coefficients.
Indeed, $|\xi x/\sqrt v|\leq1/2$ eventually. Differentiation of the
rational exponent and logarithm gives fixed polynomial bounds in
$|x|$, with denominators uniformly separated from zero. The exponential
modulus is at most $\exp(C\delta |x|^3)$, and
$C\delta X_n\leq d/2$ eventually. The complex power and its fixed
number of derivatives are bounded in the same way by polynomials.
Finite product and chain rules prove~\eqref{ncy:eq:derivbound}.

Expand~\eqref{ncy:eq:H} through degree $2S+1$. The integral Taylor
remainder, multiplied by $e^{-x^2/\sqrt v}$ and integrated, is at most
\begin{align}\label{ncy:eq:integratedrem}
 C\delta^{2S+2}(1+|\ell|)^k
       \int_{\mathbb R}P_{j,k,2S+2}(|x|)e^{-dx^2/2}\dd x
  =O\bigl(n^{-(S+1)/2}(1+|\ell|)^k\bigr).
\end{align}
This uses Gaussian moments, not the bound $|x|\leq X_n$ for the
polynomial factor, so it introduces no extra powers of $\log n$.
The identity $H_{j,k}(-\delta,-x)=H_{j,k}(\delta,x)$ implies that
the coefficient of $\delta^m$ has parity $m$ in $x$. All odd
coefficients integrate to zero on the symmetric centered segment,
also for complex $u$ because its center is the actual saddle $s$.
The remaining even coefficients can be integrated over $\mathbb R$:
the tails beyond $X_n\asymp\log n$ are smaller than every fixed
inverse power of $n$.

The moments used are
\begin{equation}\label{ncy:eq:moments}
 \int_{\mathbb R}x^{2r}e^{-x^2/\sqrt v}\dd x
  =\sqrt\pi\frac{(2r-1)!!}{2^r}v^{(2r+1)/4},\qquad
 \int_{\mathbb R}x^{2r+1}e^{-x^2/\sqrt v}\dd x=0,
\end{equation}
with $(-1)!!=1$. They hold first for real positive $v$ by substitution
and then throughout the disc by holomorphy of dominated Gaussian
integrals. This use of continuation concerns the Gaussian integral
only. Polynomial dependence on $\ell$ and holomorphic dependence on
$u$ are preserved by the finite Taylor and moment operations.
\end{proof}

\begin{proof}[Proof of Theorem~\ref{ncy:thm:main}]
For a monomial in $Q_j$, apply Lemma~\ref{ncy:lem:gaussian} with $S=J-j$.
Its $s^j$ contributes $n^{-j/2}$, and its logarithmic degree is at most
$2j$. Its total relative error is therefore
\[
 O\bigl(n^{-(J+1)/2}(1+\log n)^{2j}\bigr).
\]
The remainder in~\eqref{ncy:eq:amplitude}, integrated absolutely with the
local Gaussian, contributes
$O(n^{-(J+1)/2}(1+\log n)^{2J+2})$ relative to the leading scale.
Equation~\eqref{ncy:eq:loss}, or~\eqref{ncy:eq:exactphase} with its small
correction, supplies a fixed Gaussian majorant; $t^{-v}/s^{-v}$ is
bounded on the segment. Equation~\eqref{ncy:eq:off} disposes of the
rest of the circle.

The Jacobian and leading Gaussian integral give precisely
\begin{align*}
 \frac{n^{-3/4}}{2\pi}e^{u(cA-1)}s^{-v}e^{2\lambda}
           \sqrt\pi v^{1/4}
       =C(u)n^{v/2-3/4}e^{2\sqrt{vn}}.
\end{align*}
Every even correction coming from $Q_j$ has logarithmic degree at
most $2j$. At total order $n^{-r/2}$, $j\leq r$, proving the degree
bound $2r$ and the coefficient holomorphy.

Finally, let $L_n(u)=C(u)n^{cu/2-3/4}e^{2\sqrt{cun}}$. It is
holomorphic and nonzero on a small fixed open disc. Since $f_n(u)$ is
a polynomial, the exact expression
\begin{equation}\label{ncy:eq:residual}
 E_{n,J}(u)=\frac{f_n(u)}{L_n(u)}-1
              -\sum_{j=1}^Jn^{-j/2}R_j(\log n;u)
\end{equation}
is holomorphic there. The proof gives its bound on every sufficiently
small closed disc. Choose two concentric discs whose closures are
inside the open disc; Cauchy's derivative estimates transfer the bound
on the larger disc to all fixed derivatives on the smaller disc,
without changing its powers of $n$ or $\log n$. The contour radius
$a=\Re\sqrt{cu/n}$ need not vary holomorphically: it is a device for
proving a bound on the already holomorphic expression~\eqref{ncy:eq:residual}.
\end{proof}

\section{Explicit corrections and a fixed order calculator}\label{ncy:sec:corrections}
Retain $v=cu$, $\ell=(\log v-L)/2$ and $Q,Z,D$ from
\eqref{ncy:eq:QZ}. The first two correction polynomials are
\begin{equation}\label{ncy:eq:R1}
 R_1(L;u)=v^{3/2}Q(\ell)
              -\frac{4(v-1)^2-1}{16\sqrt v}
\end{equation}
and
\begin{align}\label{ncy:eq:R2}
 R_2(L;u)={}&\frac{[4(v-1)^2-1][4(v-1)^2-9]}{512v}\notag\\
 &+v\left\{\frac{1-4(v-2)^2}{16}Q(\ell)
              +\frac{v-2}{2}(\ell+D)-\frac14\right\}\notag\\
 &+v^2Z(\ell)+\frac{v^3}{2}Q(\ell)^2.
\end{align}
In particular the coefficient of $L^2$ in $R_1$ is $v^{3/2}/8$.

Here is a direct Gaussian check of the first correction. The
unperturbed amplitude $t^{-v}$ has order-$\delta^2$ integrand
\[
 -\frac{v+1}{2}x^2+\frac{v+1}{v^{3/2}}x^4-\frac1{2v^2}x^6.
\]
The normalized Gaussian moments are $\sqrt v/2$, $3v/4$, and
$15v^{3/2}/8$, giving $-[4(v-1)^2-1]/(16\sqrt v)$. The $tQ_1$ term
contributes $v^{3/2}Q(\ell)$, proving~\eqref{ncy:eq:R1}.

For clarity, the four parts of~\eqref{ncy:eq:R2} have separate origins.
The unperturbed fourth-order term gives its first line. The
order-$\delta^2$ correction to $tQ_1$ gives its second line. The
leading Gaussian integral of $t^2Q_2$ gives its last line, by
\eqref{ncy:eq:Q12}. Each follows from~\eqref{ncy:eq:H} and the exact moments
\eqref{ncy:eq:moments}; the companion repeats these identities symbolically.
This makes $R_2$ a finite algebraic consequence of the same local
calculation, rather than an additional asymptotic assumption.

For general fixed order, equations~\eqref{ncy:eq:polylog},
\eqref{ncy:eq:Qj}, and~\eqref{ncy:eq:H} are an algorithm using only finitely
many coefficients and Gaussian moments. An equivalent compact
formal calculator for a monomial $t^{-\beta}$ is
\begin{align}\label{ncy:eq:calculator}
 \mathcal B_\beta(v,n)={}&
 \left(\frac nv\right)^{(\beta-1)/2}
 \frac{e^{2\sqrt{vn}}}{\sqrt{4\pi\sqrt{vn}}}\notag\\
 &\times\sum_{r\geq0}
 \frac{(-1)^r\prod_{l=1}^r[4(\beta-1)^2-(2l-1)^2]}
      {r!(16\sqrt{vn})^r}.
\end{align}
For $t^{j-v}(\log t)^k$, replace the monomial by
$(-\partial_\beta)^k\mathcal B_\beta(v,n)$ at $\beta=v-j$,
holding $v$ fixed, multiply by $e^{u(cA-1)}$, and retain terms with
$j+r\leq J$. Equation~\eqref{ncy:eq:calculator} is the dominant-saddle
formal series familiar from the modified Bessel function. Its use
here is a coefficient calculator: the original full contour has not
been identified with a Bessel contour, and no convergence of the
formal series is claimed. The remainder proof is the Gaussian
argument of Section~\ref{ncy:sec:gaussian}.

\section{Refined component moments and the Gaussian limit}\label{ncy:sec:probability}
Let $K_n$ be the number of nested components in an object chosen
uniformly from the $a_n$ labeled objects. Its probability generating
function is
\[
 \EE u^{K_n}=\frac{f_n(u)}{f_n(1)}.
\]
Since $f_n/L_n\to1$ uniformly near $u=1$, the ratio is nonzero on a
fixed smaller disc for all sufficiently large $n$ and has the analytic
logarithm tending to zero there. The theorem and finite logarithmic
expansion imply
\begin{equation}\label{ncy:eq:logf}
 \log f_n(u)=\log L_n(u)+n^{-1/2}R_1(\log n;u)
          +O\bigl(n^{-1}(1+\log n)^4\bigr),
\end{equation}
with the same order for any fixed number of derivatives on a smaller
disc. This follows by applying Cauchy's estimates to the exact
normalized residual before differentiation. Thus no real-axis
asymptotic is being differentiated without justification.

\begin{theorem}[Mean and variance]\label{ncy:thm:moments}
As $n\to\infty$,
\begin{align}\label{ncy:eq:mean}
 \EE K_n&=\sqrt{cn}+\frac c2\log n+m_0
             +O\bigl(n^{-1/2}(1+\log n)^2\bigr),\\
 m_0&=-\frac34-c\log2+\frac{c\gamma}{2},\notag\\
 \Var K_n&=\frac{\sqrt{cn}}2+\frac c2\log n+v_0
             +O\bigl(n^{-1/2}(1+\log n)^2\bigr),\label{ncy:eq:variance}\\
 v_0&=-1-c\left(\log2+\frac12-\frac\gamma2\right).\notag
\end{align}
The next corrections are respectively $n^{-1/2}M_1(\log n)$ and
$n^{-1/2}V_1(\log n)$, after which both errors are
$O(n^{-1}(1+\log n)^4)$, where
\begin{equation}\label{ncy:eq:M1V1}
 M_1(L)=\left.c\partial_v\mathcal R_1(L,v)\right|_{v=c},\qquad
 V_1(L)=\left.(c\partial_v+c^2\partial_v^2)
                      \mathcal R_1(L,v)\right|_{v=c}.
\end{equation}
Here $\mathcal R_1(L,v)$ is the right side of~\eqref{ncy:eq:R1} with
$\ell=(\log v-L)/2$, and $L$ is held fixed.
\end{theorem}
\begin{proof}
The cumulant operator for marking by $u$ is
$\mathcal D=u\partial_u=v\partial_v$. Consequently the mean and
variance are $\mathcal D\log f_n|_{u=1}$ and
$\mathcal D^2\log f_n|_{u=1}$. The logarithm of the leading term is
\begin{equation}\label{ncy:eq:logleading}
 \log L_n(u)=2\sqrt{vn}+\left(\frac v2-\frac34\right)\log n
       +u(cA-1)+\left(\frac14-\frac v2\right)\log v
       -\log(2\sqrt\pi).
\end{equation}
Differentiating this explicit holomorphic function and using
$\log c=-\gamma$ gives~\eqref{ncy:eq:mean}--\eqref{ncy:eq:variance}.
Differentiating the next term in~\eqref{ncy:eq:logf} gives
\eqref{ncy:eq:M1V1}; its residual is controlled as just explained.
\end{proof}

An entirely explicit form of these first refined polynomials is also
useful. Write $q=Q(\ell_c)$, $q'=\ell_c+D$, and
$\ell_c=(\log c-L)/2$. Then
\begin{align}\label{ncy:eq:explicitM1}
 M_1(L)={}&c^{3/2}\left(\frac32q+\frac12q'\right)
       -\frac38c^{3/2}+\frac14c^{1/2}+\frac3{32}c^{-1/2},\\
 V_1(L)={}&c^{3/2}\left(\frac94q+\frac32q'+\frac14\right)
       -\frac9{16}c^{3/2}+\frac18c^{1/2}-\frac3{64}c^{-1/2}.
       \label{ncy:eq:explicitV1}
\end{align}
These follow by writing the nonlogarithmic term in~\eqref{ncy:eq:R1}
as $-v^{3/2}/4+v^{1/2}/2-3v^{-1/2}/16$ and using
$v\partial_v\ell=1/2$.

In particular,
\begin{equation}\label{ncy:eq:meanlimit}
 \lim_{n\to\infty}\frac{\EE K_n}{\sqrt n}
       =e^{-\gamma/2}\approx0.749306001288449.
\end{equation}
The decimal is uncertified. Riedel's numerical estimate
$0.7719351399$~\cite[p. 3]{riedel} is a finite-data estimate of the
constant, accompanied there by caution about the fitted exponent.
The positive $\frac c2\log n$ term explains a slow approach to the
limit. The exact square-root limit belongs to the pre-existing
expansive-assembly regime; the displayed lower-order corrections are
specific to these component weights.

\begin{corollary}[Credited Gaussian regime]\label{ncy:cor:clt}
With $\sigma_n^2=\sqrt{cn}/2$,
\begin{equation}\label{ncy:eq:clt}
 \frac{K_n-\EE K_n}{\sigma_n}\ \Longrightarrow\ N(0,1).
\end{equation}
The same statement holds with $\sigma_n^2$ replaced by $\Var K_n$.
\end{corollary}
\begin{proof}
Use the analytic logarithm of $f_n(e^t)/f_n(1)$ in a fixed complex
neighborhood of $t=0$. By~\eqref{ncy:eq:logf} and Cauchy's estimates,
its third derivative is $O(\sqrt n+\log n)$ there. Its second derivative
at $0$ is $\sigma_n^2+O(\log n)$. For fixed real $\theta$, substitute
$t=i\theta/\sigma_n$ and subtract $t\EE K_n$. Taylor's theorem gives
\[
 \log\EE\exp\left(i\theta\frac{K_n-\EE K_n}{\sigma_n}\right)
 =-\frac{\theta^2}{2}+o(1),
\]
with cubic error $O(n^{-1/4})$ and second-derivative replacement error
$O(n^{-1/2}\log n)$. The continuity theorem for characteristic
functions proves~\eqref{ncy:eq:clt}. Equation~\eqref{ncy:eq:variance} and
Slutsky's theorem justify the variance normalization.
\end{proof}

The proof recovers a known generic limit-law phenomenon with the
normalization needed here. It is not a new universality assertion.
No local limit theorem or quantitative distribution-distance bound
is inferred from this fixed-frequency argument.

\begin{remark}[{[write]} Riedel's conjecture, and what his fitted constants missed, 6~October 2026]\label{ncy:rem:riedel}
\emph{The source.} Riedel's note~\cite{riedel} has a paragraph ``The
conjecture'' on p.~3 of~6, the same in the revised version (page footers
dated 24.01.26) and in the original linked from A392471 (footers dated
13.01.26). With $X$ the number of components of a uniformly random nested
cycle partition on $n$ nodes, and $\mathrm E[X]$ computed from his
recurrences, the boxed conjecture is
\[
 \mathrm E[X]\sim N\sqrt n ,
\]
where the constant $N$ is left to be determined (readers able to contribute
are asked to contact him). He adds that the asymptotics might instead
involve a power $n^M$ with $M$ ``close to but not equal to $1/2$'', and
reports that the numerical data suggest $M\ge0.4704567259$ and
$N\approx0.7719351399$. The note gives neither the range of $n$ nor the
fitting method. Neither OEIS entry contains the conjecture
(Section~\ref{ncy:sec:provenance}). His $X$ is $K_n$: his EGF is
$Q(z,u)=\exp\bigl(u\sum_{k\ge1}Q_k(z)\bigr)$ with
$\sum_{k\ge1}Q_k(z)=P(z)-1$ (his p.~2), which is $F(z,u)$
of~\eqref{ncy:eq:F}.

\emph{(a) What is proved.} By~\eqref{ncy:eq:meanlimit} the boxed conjecture
holds, with $N=e^{-\gamma/2}=0.7493060\ldots$; this answers his question
for~$N$. The alternative is excluded: for $N>0$, $\EE K_n\sim Nn^M$ holds
only for $M=1/2$, since $\EE K_n/n^M=(\EE K_n/\sqrt n)\,n^{1/2-M}$ tends to
$+\infty$ for $M<1/2$ and to~$0$ for $M>1/2$. The inequality
$M\ge0.4704567259$, as printed, is satisfied by $M=1/2$; the value
$N\approx0.7719351399$ is not the constant: it exceeds $e^{-\gamma/2}$ by
$0.0226$. The source proves all of this in Theorem~\ref{ncy:thm:moments} but
calls Riedel's number a finite-data estimate and does not say that the
conjecture is settled; it is.

\emph{(b) What the fits missed} [write]. Put $E_n=\EE K_n$ and
$m_0=-0.977132\ldots$ as in~\eqref{ncy:eq:mean}. Then
\begin{equation*}
 \frac{E_n}{\sqrt n}=e^{-\gamma/2}+\frac{\frac c2\log n+m_0}{\sqrt n}
   +O\!\left(\frac{(1+\log n)^2}{n}\right),
\end{equation*}
and the numerator $\frac c2\log n+m_0$ is positive exactly when
$n>e^{-2m_0/c}=32.48\ldots$; so $E_n/\sqrt n>e^{-\gamma/2}$ for all large
$n$ and approaches it only at the rate $(\log n)/\sqrt n$. For every fixed
integer $\varrho\ge2$ the two-point exponent satisfies
\begin{equation*}
 M_\varrho(n):=\frac{\log(E_{\varrho n}/E_n)}{\log\varrho}
  =\frac12-\bigl(1-\varrho^{-1/2}\bigr)\frac{c\log n}{2\sqrt{cn}\,\log\varrho}
   +O\bigl(n^{-1/2}\bigr),
\end{equation*}
so it lies below $1/2$ for all large $n$ and tends to $1/2$ at the same slow
rate. \emph{Proof.} By~\eqref{ncy:eq:mean},
$E_n=\sqrt{cn}\,(1+\delta_n)$ with
$\delta_n=(\frac c2\log n+m_0)/\sqrt{cn}+O(n^{-1}(1+\log n)^2)\to0$. Since
$\log(1+x)=x+O(x^2)$,
$\log(E_{\varrho n}/E_n)=\frac12\log\varrho+\delta_{\varrho n}-\delta_n
+O(n^{-1}(1+\log n)^2)$, and
\[
\begin{aligned}
 \delta_{\varrho n}-\delta_n
 &=\frac{\varrho^{-1/2}\bigl(\frac c2\log n+\frac c2\log\varrho+m_0\bigr)
   -\bigl(\frac c2\log n+m_0\bigr)}{\sqrt{cn}}+O\!\left(\frac{(1+\log n)^2}{n}\right)\\
 &=-\bigl(1-\varrho^{-1/2}\bigr)\frac{c\log n}{2\sqrt{cn}}+O\bigl(n^{-1/2}\bigr).
\end{aligned}
\]
Dividing by $\log\varrho$ gives the claim. $\square$ A fitted exponent below
$1/2$ and a fitted constant above $e^{-\gamma/2}$ are therefore what the
proved law produces at every large finite $n$: the additive correction
$\frac c2\log n+m_0$ in~\eqref{ncy:eq:mean}, which is not a power, is read
by a fit as a change of exponent and of constant. Values (double-precision
evaluation of $E_n$ from recurrences with positive terms, agreeing with the
exact rational $E_{100}$ and with a 256-bit fixed-point evaluation at
$n=4500$):
\[
\begin{array}{r|ccccc}
 n & 100 & 1000 & 4500 & 10000 & 20000\\\hline
 E_n/\sqrt n & 0.821954 & 0.786471 & 0.771935 & 0.766422 & 0.762634\\
 M_2(n) & 0.47781 & 0.48628 & 0.49089 & 0.49285 & 0.49428
\end{array}
\]
For $1\le n\le40000$ every computed ratio $E_n/\sqrt n$ exceeds
$e^{-\gamma/2}$, and the ratios decrease from $n=4$ on.

\emph{(c) Where Riedel's numbers may come from.}
$E_{4500}/\sqrt{4500}=0.77193514015$ agrees with his $N\approx0.7719351399$
to $2.5\times10^{-10}$, and the ratio first falls below his value at
$n=4501$. This suggests, but the note does not say, that his $N$ is the
finite ratio at $n=4500$. The write could not identify the computation behind
$M\ge0.4704567259$: the one-step log-log slopes
$\log(E_n/E_{n-1})/\log(n/(n-1))$ for $n\le40000$, the two-point exponents
$\log(E_{4500}/E_m)/\log(4500/m)$ for $m<4500$, and least-squares fits of
$\log E_n$ against $\log n$ over the windows $[a,4500]$, $1\le a<60$, and a
few shorter windows were tried, and none reproduces it to its ten digits. Recorded only: nothing was submitted to the OEIS and
Riedel was not contacted.
\end{remark}

\section{Asymptotic inversion with safe integer rounding}\label{ncy:sec:inverse}
For sufficiently large real $y$, define
\[
 N(y)=\min\{n\geq0:a_n\geq y\}.
\]
The sequence is strictly increasing for $n\geq1$. Adding label $n+1$
as a new singleton component is an injection; a single cycle on all
$n+1$ labels is outside its image. The equality $a_0=a_1=1$ causes
no problem for large $y$.

Put $\beta=c/2-3/4$. For fixed $J\geq0$, define at sufficiently large
real $x$
\begin{align}\label{ncy:eq:G}
 G_J(x)={}&\log\Gamma(x+1)+2\sqrt{cx}+\beta\log x+\log C\notag\\
 &+\log\left(1+\sum_{j=1}^Jx^{-j/2}R_j(\log x;1)\right).
\end{align}
The last logarithm is real and well defined eventually, and
$G_J'(x)\sim\log x$. At integer $n$, Theorem~\ref{ncy:thm:main} implies
\begin{equation}\label{ncy:eq:logerror}
 \log a_n=G_J(n)+O(n^{-q}(\log n)^p),\qquad
 q=(J+1)/2,\quad p=2J+2.
\end{equation}
For large $y$, let $x_J$ be the unique sufficiently large real solution
of $G_J(x_J)=\log y$.

\begin{theorem}[Two-ceiling enclosure]\label{ncy:thm:inverse}
For each fixed $J$ there exist $M_J>0$ and a threshold such that, with
\[
 \epsilon_J=M_Jx_J^{-q}(\log x_J)^{p-1},
\]
one has
\begin{equation}\label{ncy:eq:inverse}
 \lceil x_J-\epsilon_J\rceil\ \leq N(y)\ \leq
                       \lceil x_J+\epsilon_J\rceil.
\end{equation}
If $\dist(x_J,\mathbb Z)>\epsilon_J$, the two ceilings coincide and
$N(y)=\lceil x_J\rceil$. No explicit certified $M_J$ or onset threshold
is supplied.
\end{theorem}
\begin{proof}
Bound the integer error in~\eqref{ncy:eq:logerror} by
$Kn^{-q}(\log n)^p$. For $x=x_J$ sufficiently large,
$\epsilon_J<1$, $G_J'\geq(\log x)/2$ on $[x-2,x+2]$, and the
error at integers in that interval is at most $2Kx^{-q}(\log x)^p$.
Choose $M_J>4K$. Then $n_+=\lceil x+\epsilon_J\rceil$ satisfies
\[
 G_J(n_+)-G_J(x)\geq\frac{\log x}{2}\epsilon_J
                       >2Kx^{-q}(\log x)^p,
\]
so $\log a_{n_+}>\log y$. Similarly
$n_-=\lceil x-\epsilon_J\rceil-1<x-\epsilon_J$ satisfies
$\log a_{n_-}<\log y$. Monotonicity gives~\eqref{ncy:eq:inverse}.
Enlarging the threshold ensures all comparisons used above.
\end{proof}

A single ceiling of a finite continuous asymptotic expansion is not
valid unconditionally: a target may lie arbitrarily close to an exact
integer threshold. One must use the two-ceiling enclosure, prove the
stated separation, or compare exact counts.

For a convenient starting approximation let $L=\log y$ and use the
positive real branch of Lambert's $W$ function:
\begin{equation}\label{ncy:eq:W}
 x_0^*=\frac{L}{W(L/e)}.
\end{equation}
It solves $x(\log x-1)=L$, and $N(y)\sim x_0^*$. With
$\alpha=c/2-1/4$ and $K_0=\log(C\sqrt{2\pi})$, Stirling's formula gives
\begin{align}\label{ncy:eq:newton}
 x_*&=x_0^*-
 \frac{2\sqrt{cx_0^*}+\alpha\log x_0^*+K_0}{\log x_0^*},\\
 x_J&=x_*+O\bigl((\log x_0^*)^{-2}\bigr).\notag
\end{align}
For justification, the displacement from $x_0^*$ is
$O(\sqrt{x_0^*}/\log x_0^*)$. Substituting~\eqref{ncy:eq:newton} into
$G_J-L$, the uncanceled derivative correction is $O(1/\log x_0^*)$;
the quadratic curvature term is $O((\log x_0^*)^{-2})$ and the
omitted Stirling and fixed-order corrections are smaller. Dividing
by $G_J'\sim\log x_0^*$ proves the displayed error. Newton iteration
or formal reversion of~\eqref{ncy:eq:G} gives further continuous orders.
Every integer conclusion still needs~\eqref{ncy:eq:inverse} or an exact
comparison. The companion's bounded exact inverse compares integer
counts directly and does not use an uncertified asymptotic rounding
rule.

\begin{remark}[{[write]} The inversion against the transseries volume, 6~October 2026]\label{ncy:rem:transseries}
The volume is
\path{Analysis/Transseries/docs/series-and-transseries/Transseries_And_Inversion/};
its letters $\kappa$, $d$, $v$, $X$, $L$ are written here with the subscript
``vol'' where they would collide with this report's.
\begin{enumerate}
\item[(a)] \emph{The starting value $x_0^\ast$ of~\eqref{ncy:eq:W} is an
 instance}, verbatim, of \file{p0:prop:factorial-core} with
 $(\kappa_{\mathrm{vol}},d_{\mathrm{vol}},L_{\mathrm{vol}}):=(1,-1,L)$,
 $L=\log y>0$: the phase is $X(\log X-1)$, the solution is
 $v_{\mathrm{vol}}=W_0(L/e)$ and
 $X_{\mathrm{vol}}=e^{v_{\mathrm{vol}}+1}=L/v_{\mathrm{vol}}=x_0^\ast$ (using
 $v_{\mathrm{vol}}e^{v_{\mathrm{vol}}}=L/e$), and the side condition
 $\log X_{\mathrm{vol}}-1>0$ holds because $v_{\mathrm{vol}}>0$. The
 ``positive real branch'' of the source is $W_0$.
\item[(b)] \emph{$N(y)$ is the staircase} $N_\ast(y)$ of
 \file{p0:def:three-inverses} for the sequence $(a_n)$ with $n_1=1$, for
 every $y>1$: the sequence is strictly increasing for $n\ge1$ (the injection
 argument at the start of this section), and for $y>1=a_0=a_1$ the minima over
 $n\ge0$ and over $n\ge1$ agree.
\item[(c)] \emph{Theorem~\ref{ncy:thm:inverse} is not an instance} of
 \file{p0:thm:staircase}, only an analogue of its item~(2). Its proof compares
 the model $G_J$ with $\log a_n$ at the integers and then uses monotonicity
 directly; no admissible interpolation of $(a_n)$ appears, and $x_J$ is the
 root of the model, not the inverse $\nu_{\mathcal A}(y)$ of an
 interpolation, so the hypothesis $|x_\ast-\nu(y)|\le\eta$ of item~(2) is not
 established. The conclusion has the two-ceiling form of item~(2), and the
 paragraph after the proof describes the failure mode stated there.
 The Newton correction~\eqref{ncy:eq:newton} is proved directly; no
 comparison with the volume's perturbation series is claimed.
\end{enumerate}
No novelty is claimed for these inversions.
\end{remark}

\section{Exact recurrences and reproducible checks}\label{ncy:sec:computation}
The infinite products in the definitions reduce to finite exact
computations at every fixed size. The companion uses a deliberately
bounded range, not an unbounded enumeration promise.

Let $B_n^{(m)}$ count labeled permutations of $n$ labels having at
most one cycle of each length in $\{1,\ldots,m\}$. With
$B_0^{(0)}=1$ and all other entries zero,
\begin{equation}\label{ncy:eq:integercycle}
 B_n^{(m)}=B_n^{(m-1)}+
 \begin{cases}
  \binom nm(m-1)!B_{n-m}^{(m-1)},& n\geq m,\\
  0,&n<m.
 \end{cases}
\end{equation}
The cycle's labels are chosen first and its cyclic order next; using
the previous row enforces at most one such cycle. Thus $B_n=B_n^{(n)}$.
The component containing label $1$ gives the partial Bell recurrence
\begin{equation}\label{ncy:eq:bell}
 T(n,k)=\sum_{m=1}^n\binom{n-1}{m-1}B_mT(n-m,k-1),
 \qquad n,k\geq1.
\end{equation}
All arithmetic in these recurrences is integer arithmetic. Their
first five nonempty rows and row sums are
\[
\begin{array}{c|l|r}
 n&T(n,1),\ldots,T(n,n)&a_n\\\hline
 1&1&1\\
 2&1,1&2\\
 3&5,3,1&9\\
 4&14,23,6,1&44\\
 5&74,120,65,10,1&270
\end{array}
\]

An independent rational-series check follows from
\[
 \log P(z)=\sum_{r\geq1}\ell_rz^r,\qquad
 \ell_r=\sum_{m\mid r}\frac{(-1)^{r/m+1}}{(r/m)m^{r/m}}.
\]
Formal differentiation gives
\begin{equation}\label{ncy:eq:rational}
 np_n=\sum_{r=1}^n r\ell_rp_{n-r},\qquad
 nf_n(1)=\sum_{r=1}^n rp_rf_{n-r}(1),\qquad p_0=f_0(1)=1.
\end{equation}
A separate power-series multiplication checks
$T(n,k)=n![z^n](P-1)^k/k!$ directly. At small sizes, literal
permutation cycle decomposition followed by enumeration of admissible
groupings of cycles gives an additional combinatorial check.

The frozen OEIS fixtures contain the displayed numeric terms with
public source URLs and retrieval metadata: 23 values of A007838,
22 values of A308338, and 55 entries of A392471. They are independent
expected values, not values generated during a verification run.
The companion checks those prefixes, complete counts and triangle
through $n=100$, independent rational count recurrences through
$n=100$, power-series triangle identities through $n=30$, and literal
enumeration through $n=7$. Exact rational mean and variance samples
come directly from the triangle. Optional SymPy checks verify the
finite symbolic correction identities. See the package README for
the commands and dependency versions used.

\textbf{Evidence levels.} The fixtures, finite recurrences, exact
integer comparisons, rational moments, and symbolic identities are
exact. They are finite checks and do not prove an asymptotic theorem.
The fixed-order remainder, complex uniformity, all-circle suppression,
moment expansion, CLT, and inverse enclosure are proved in the text
using the credited component estimate. Their implied constants and
onset thresholds have not been made effective here. All decimal
constants or other numerical asymptotic diagnostics are uncertified;
high arithmetic precision alone would not certify their errors.
The package supplies no interval enclosure for the transcendental
constant $C$ and no empirical fitted constant as a substitute for its
exact definition.

\begin{writenote}[The package here, and the reruns, 6~October 2026]
``The companion'' and ``the package README'' above are the delivered ones. In
this report the companion programs are
\file{code/companion-exact_nested.py}, \file{code/companion-verify.py},
\file{code/companion-test_exact_nested.py} and
\file{code/companion-symbolic_checks.py} (delivered in \file{companion/});
the frozen fixtures and outputs are in \file{data/}; the companion's README
and provenance note are \file{companion-README.md} and
\file{companion-PROVENANCE.md}; the release tools \file{build_pdf.py},
\file{make_zip.py}, \file{release_tools.py} and \file{test_release.py} are in
\file{code/}; \file{SOURCE_PROVENANCE.json} and
\file{verification_receipt.json} are in \file{data/}; the checksum manifest
is not shipped. \file{verify.py} reads \file{../data/published_fixtures.json}
relative to its own folder, so the programs run only on a copy with the
delivered names (the report README gives two routes). On such copies at the
write (Windows, Python~3.14.4, SymPy~1.14.0): from a fresh extraction of the
archive, \texttt{verify.py}, \texttt{verify.py --include-data},
\texttt{test\_exact\_nested.py} and \texttt{symbolic\_checks.py} reproduce
\file{full_verification.json}, \file{exact_data_100.json},
\file{tests.normal.json} and \file{symbolic.normal.json} once Windows line
endings are converted; from the shipped files, the verifier and the
regression tests reproduce their outputs the same way. The release tests need
POSIX descriptor support and were not run.
\end{writenote}

\section{Source coverage and attribution}\label{ncy:sec:sources}
The source review covered the three OEIS entries and Riedel's revised
January 2026 note, the 2006 hybrid asymptotics paper, and the later
Erlihson--Granovsky expansive-assembly paper. It also checked the
connection between the stated component weights and the hypotheses
of that general assembly theory. This is bounded source coverage,
not a worldwide historical-priority certificate.

The attribution is deliberately separated by mathematical role.
Gutkovskiy's 2019 OEIS entry supplies A308338's defining EGF. Riedel
supplies the nested-cycle interpretation and A392471. Greene and
Knuth are credited for~\eqref{ncy:eq:prior} by the inspected hybrid paper;
the original book was not independently inspected. The hybrid paper
itself develops the full component expansion, with root-of-unity
periodic coefficient terms. Those component results do not by
themselves establish the marked outer-exponential expansion proved
here.

The inspected Erlihson--Granovsky preprint is the 19 June 2007
arXiv version of the paper published in October 2008. Theorem and
corollary numbers agree with the journal version, while some equation
numbers differ. The reference to the 2004 total-component CLT is
attribution verified in the later work, not a direct reading of the
2004 paper. That original paper was not inspected for this report.

The variance in~\eqref{ncy:eq:variance} is derived independently from
\eqref{ncy:eq:logleading}. There is a normalization inconsistency in the
later source's displayed general-$C$ covariance formula (preprint
(4.54), published (4.18)): with its definition
$b_r(u)=C\Gamma(r+1,u)$, the printed conditioning denominator is
$\Gamma(p+2)$. The hypotheses allow every $C>0$. At the interior choice
$p=1$, $C=3$, $u=1$, literal substitution gives
$3/e-18/e^2<0$, which cannot be a variance.
This display is therefore not used to certify our constant. The
paper remains the cited source for the general regime and its explicit
historical attribution.

Within the checked sources, the explicit count prefactor, fixed-order
complex-marked outer expansion, logarithmic moment refinements, and
rounding-qualified inverse are the sequence-specific development of
this report. No claim is made that no other work contains overlapping
results. There is no external publication or amendment of an OEIS
entry associated with this document.

\begin{remark}[{[write]} The covariance display of Erlihson and Granovsky, corrected, 6~October 2026]\label{ncy:rem:eg}
In this remark $p$, $C$, $u$, $q$, $N$, $r_N$, $\delta$, $f_j$, $b_r$,
$\nu$, $S_j$, $Z_N$ are the letters of~\cite{eg2008}: the exponent and
constant in $a_k\sim Ck^{p-1}$ ($k\to\infty$; every $p>0$, $C>0$ is allowed),
a scaled component size, the number of strata, the total size, the scaling
$r_N=h^{-1}N^{1/(p+1)}$ with $h=(C\Gamma(p+1))^{1/(p+1)}$, $\delta=1/r_N$,
and so on; they are not $p_m$, the amplitude~$C$, the marking variable~$u$
or the $q$ of Section~\ref{ncy:sec:inverse}.

\emph{The display.} With $b_r(u)=C\Gamma(r+1,u)$ (preprint~(4.53),
published~(4.17)), Theorem~4.3(i) states that the scaled counts
$(\hat\nu_1,\dots,\hat\nu_q)$ of components of size at least $u_jr_N$
converge to a centred Gaussian vector with covariance
\[
 e_{mk}=b_{p-1}(u_s)-\frac{b_p(u_m)b_p(u_k)}{\Gamma(p+2)},\qquad s=\max(k,m)
\]
(preprint~(4.54); published~(4.18) on p.~927, identical, read by the write),
and Corollary~4.5 ($q=1$, $u_1=u$) gives $\hat\nu(u)\Rightarrow N(0,e_{11}(u))$
``for all $u\ge0$'', the case $u=0$ being the total number of components.
The scaling of~$\hat\nu_j$ (preprint~(4.45), published~(4.9)) multiplies
$\nu_j-S_j$ by $(f_0(p-1)/S_0)^{1/2}\sim r_N^{-p/2}$.

\emph{(a) The display fails.} At $p=1$, $C=3$ (for instance $a_k=3$ for all
$k$), $q=1$, $u_1=1$: $b_0(1)=3\Gamma(1,1)=3/e$, $b_1(1)=3\Gamma(2,1)=6/e$
and $\Gamma(3)=2$, so
\[
 e_{11}=\frac3e-\frac{36}{2e^2}=\frac3e-\frac{18}{e^2}=-1.33240\ldots<0,
\]
which is not the variance of any random variable. More generally, at $u=0$
the display gives
$e_{11}(0)=C\Gamma(p)-C^2\Gamma(p+1)^2/\Gamma(p+2)=C\Gamma(p)\bigl(1-Cp/(p+1)\bigr)$,
negative whenever $C>(p+1)/p$. The source found the first value (Section~\ref{ncy:sec:sources}).

\emph{(b) The correction.} The denominator should be
$C\Gamma(p+2)=b_{p+1}(0)$:
\[
 e^{\mathrm{corr}}_{mk}=b_{p-1}(u_s)-\frac{b_p(u_m)b_p(u_k)}{C\,\Gamma(p+2)},
\]
which coincides with the printed display exactly when $C=1$. This is what the
conditioning computation gives. In the Boltzmann model at $\delta=1/r_N$ the
counts $n_k$ are independent Poisson variables of means $a_ke^{-\delta k}$;
with $\nu_j=\sum_{k\ge u_jr_N}n_k$ and the total mass $Z=\sum_kkn_k$, Riemann
sums give
$r_N^{-p}\operatorname{Cov}(\nu_m,\nu_k)\to C\Gamma(p,u_s)=b_{p-1}(u_s)$,
$r_N^{-p-1}\operatorname{Cov}(\nu_m,Z)\to C\Gamma(p+1,u_m)=b_p(u_m)$ and
$r_N^{-p-2}\operatorname{Var}Z\to C\Gamma(p+2)=b_{p+1}(0)$, and the
covariance of a Gaussian vector conditioned on $Z$ is
$\operatorname{Cov}(\nu_m,\nu_k)-\operatorname{Cov}(\nu_m,Z)\operatorname{Cov}(\nu_k,Z)/\operatorname{Var}Z$.
The corrected display is a variance: $e^{\mathrm{corr}}_{11}(u)
=C\bigl(\Gamma(p,u)-\Gamma(p+1,u)^2/\Gamma(p+2)\bigr)>0$, because the
Cauchy--Schwarz inequality for $x^{(p-1)/2}e^{-x/2}$ and
$x^{(p+1)/2}e^{-x/2}$ on $(u,\infty)$ gives
$\Gamma(p+1,u)^2<\Gamma(p,u)\Gamma(p+2,u)\le\Gamma(p,u)\Gamma(p+2)$. In the
example of~(a) it is $3/e-6/e^2=0.29163\ldots$. The write checked these
limits and the conditioning formula, not the remaining steps of Section~5
of~\cite{eg2008}.

\emph{(c) Where the factor $C$ is lost.} The published (5.62) on p.~936
(preprint~(5.121)) reads
$\operatorname{Var}Z_N=\sum_kk^2a_ke^{-\delta k}\sim\Gamma(p+2)\delta^{-(p+2)}$;
under $a_k\sim Ck^{p-1}$ the right side is $C\Gamma(p+2)\delta^{-(p+2)}$.
The same $\Gamma(p+2)$ without~$C$ enters $T_q$ in (5.69) (preprint~(5.128)),
while the proof of (5.72) on p.~938 invokes
$\sum_jf_j(p+1)=C\Gamma(p+2)$. The covariance~(4.13) of Theorem~4.1(ii)
(preprint~(4.49); published p.~926) carries the same denominator and fails
in the same way: at $q=1$, $p=1$, $C=3$, $u_1=10$ its
entry $f_0(0)-f_0(1)^2/\Gamma(3)$ equals $-1.4956\ldots$ (corrected:
$1.5014\ldots$). The write did not trace which step first drops~$C$.

\emph{(d) The display against this report.} The model here has weights
$p_m\to c$ (Section~\ref{ncy:sec:exact}), that is $p=1$ and $C=c$ in the
letters of~\cite{eg2008}, with $N=n$ and $r_N=(n/c)^{1/2}$. Read at $u=0$, as
their Corollary~4.5 is read there, with the normalization $r_N^{-1/2}$ (the
limit of their scaling factor for $u_1>0$), the printed display gives
$(K_n-S)/r_N^{1/2}\Rightarrow N\bigl(0,c(1-c/2)\bigr)=N(0,0.40384\ldots)$ for
a deterministic centring~$S$. Corollary~\ref{ncy:cor:clt}
and~\eqref{ncy:eq:variance}, proved here without their display, give
$(K_n-\EE K_n)/r_N^{1/2}\Rightarrow N(0,c/2)=N(0,0.28072\ldots)$, because
$\Var K_n/r_N\to c/2$. The two normalized variables differ by a deterministic
shift~$t_n$, and if $Y_n\Rightarrow N(0,\alpha)$ and $Y_n+t_n\Rightarrow
N(0,\alpha')$ then $(t_n)$ is bounded and each of its limit points~$t$ gives
$N(t,\alpha)=N(0,\alpha')$, hence $\alpha=\alpha'$. So the printed display
contradicts Corollary~\ref{ncy:cor:clt} in this model, and the corrected one,
$e^{\mathrm{corr}}_{11}(0)=c-c^2/(2c)=c/2$, agrees with it.

The source used the display for nothing and derived~\eqref{ncy:eq:variance}
independently; no report in the repository uses it. Recorded only: nothing
was sent to the authors.
\end{remark}

\section{Further questions}\label{ncy:sec:questions}
\begin{writenote}[Added 6 October 2026, batch~108]
Following Vladimir Reshetnikov's standing rule of 4~October 2026, every
claim of the source that is not proved in full is stated here as an open
question. Questions~1--5 are the source's, as delivered; what each lacks is
what it names. Question~1: explicit constants in~\eqref{ncy:eq:prior}, in
the off-saddle bound~\eqref{ncy:eq:off} and in
Lemma~\ref{ncy:lem:gaussian}. Question~2: control of $f_n(u)$ for $u$ beyond
a small disc about~$1$. Question~3: a global analysis of the secondary arcs
beyond the all-arc bound~\eqref{ncy:eq:globalP}. Question~4: complexity
bounds and certified simplification for the calculator. Question~5: a faster
exact algorithm with an independent check. The source's non-claims are
collected in Section~\ref{ncy:sec:provenance}. The intake found no false
claim in the source and the write found none either; the two corrections
concern outside literature (Remarks~\ref{ncy:rem:riedel}
and~\ref{ncy:rem:eg}). Question~1 would also give the onset in
Remark~\ref{ncy:rem:riedel}(b): $\EE K_n/\sqrt n$ exceeds $e^{-\gamma/2}$ for
all large $n$, and for every computed $n\le40000$, but no explicit $n_0$ is
proved. Question~6 is the write's.
\end{writenote}
Several extensions require additional work beyond the proved fixed-order
statements.
\begin{enumerate}
 \item\label{ncy:q:effective} Can the component estimates and Gaussian remainders be made
 effective enough to give certified constants and usable onset
 thresholds for the count equivalent and inverse enclosure?
 \item\label{ncy:q:local} Can complex marking be controlled beyond a small disc to prove
 a lattice local limit theorem, moderate deviations, or a quantitative
 normal-approximation bound for $K_n$?
 \item\label{ncy:q:periodic} How do the periodic component corrections at secondary roots
 of unity appear in exponentially small outer-coefficient terms?
 A complete beyond-all-orders expansion would require a separate
 global analysis and is not implied by the present all-arc bound.
 \item\label{ncy:q:calculator} Can the fixed-order coefficient calculator be extended to
 higher corrections with explicit complexity bounds and certified
 symbolic simplification?
 \item\label{ncy:q:bell} Can the exact Bell triangle be computed faster at large $n$
 while retaining an independent check and exact integer arithmetic?
 \item\label{ncy:q:eg} \textup{[write]} Remark~\ref{ncy:rem:eg} corrects
 the covariance displays (4.13) and (4.18) of Erlihson and
 Granovsky~\cite{eg2008} by restoring the factor $C$ in the variance of the
 total mass, and shows that the printed displays fail. Does their proof, with
 (5.62) and (5.69) corrected, give exactly the corrected displays, and does
 any other statement of their Sections~4 and~5 depend on the missing factor?
 What is missing is a line-by-line audit of their Section~5; the write
 checked only the conditioning computation and the three limits it uses.
\end{enumerate}
These are research questions, not consequences of the finite numerical
or symbolic checks.

\begin{thebibliography}{9}
\bibitem{oeis7838}
The OEIS Foundation, \emph{A007838}, permutations with distinct cycle
lengths. \url{https://oeis.org/A007838}. Entry inspected 3 October 2026.

\bibitem{oeis308338}
I. Gutkovskiy, \emph{A308338}, OEIS entry dated 20 May 2019.
\url{https://oeis.org/A308338}. Entry inspected 3 October 2026.

\bibitem{oeis392471}
M. Riedel, \emph{A392471}, triangle counting nested-cycle assemblies
by component count. \url{https://oeis.org/A392471}.
Entry inspected 3 October 2026.

\bibitem{riedel}
M. Riedel, \emph{Nested Cycle Partitions: a conjecture}, revised 24 January 2026
author-hosted note, 6 pp.
\url{https://pnp.mathematik.uni-stuttgart.de/iadm/Riedel/papers/ncp-comp2.pdf}.
Original 13 January version linked from OEIS:
\url{https://oeis.org/A392471/a392471.pdf}.
The nested-cycle specification and component-count discussion were
inspected directly.

\bibitem{hybrid}
P. Flajolet, \'E. Fusy, X. Gourdon, D. Panario, and N. Pouyanne,
\emph{A Hybrid of Darboux's Method and Singularity Analysis in
Combinatorial Asymptotics}, Electronic Journal of Combinatorics
\textbf{13}(1) (2006), R103, 35 pp.
\url{https://arxiv.org/abs/math/0606370}.
Section 3 and Proposition 1 of the 15 June 2006 version were inspected.
The Greene--Knuth attribution is indirect through that work; the book
itself was not inspected.

\bibitem{eg2008}
M. M. Erlihson and B. L. Granovsky,
\emph{Limit shapes of Gibbs distributions on the set of integer
partitions: The expansive case}, Annales de l'Institut Henri Poincar\'e,
Probabilit\'es et Statistiques \textbf{44}(5) (2008), 915--945.
\url{https://doi.org/10.1214/07-AIHP129}.
Published text: \url{https://numdam.org/item/10.1214/07-AIHP129.pdf}.
The assembly weights, expansive scaling, Theorems 4.1 and 4.3, and
Corollary 4.5 were checked against the 2007 preprint; relevant published
equation numbering and covariance text were also checked.

\bibitem{eg2004}
M. Erlihson and B. Granovsky,
\emph{Reversible coagulation-fragmentation processes and random
combinatorial structures: asymptotics for the number of groups},
Random Structures \& Algorithms \textbf{25} (2004), 227--245.
Cited here solely through the explicit historical attribution in
\cite[Corollary 4.5 and reference 12]{eg2008}; this paper was not
inspected directly.
\end{thebibliography}
\end{document}
